\hypertarget{cpu-io-bound-system-concurrency}{%
\chapter{CPU \& I/O BOUND SYSTEM
CONCURRENCY}\label{cpu-io-bound-system-concurrency}}

\hypertarget{concurrency-paradigms-threads-vs.-processes-vs.-coroutines-asyncio}{%
\subsection{27.1 Concurrency Paradigms: Threads vs.~Processes
vs.~Coroutines
(Asyncio)}\label{concurrency-paradigms-threads-vs.-processes-vs.-coroutines-asyncio}}

Python supports three core concurrency paradigms, each targeting
different system limitations. The table below outlines their
architectural differences:

\begin{longtable}[]{@{}
  >{\raggedright\arraybackslash}p{(\columnwidth - 10\tabcolsep) * \real{0.17}}
  >{\raggedright\arraybackslash}p{(\columnwidth - 10\tabcolsep) * \real{0.17}}
  >{\raggedright\arraybackslash}p{(\columnwidth - 10\tabcolsep) * \real{0.17}}
  >{\raggedright\arraybackslash}p{(\columnwidth - 10\tabcolsep) * \real{0.17}}
  >{\raggedright\arraybackslash}p{(\columnwidth - 10\tabcolsep) * \real{0.17}}
  >{\raggedright\arraybackslash}p{(\columnwidth - 10\tabcolsep) * \real{0.17}}@{}}
\toprule
Paradigm & Execution Type & GIL Limitation & Primary Use Case & Overhead
& Communication Method \\
\midrule
\endhead
\textbf{Multithreading} (\passthrough{\lstinline!threading!}) &
Preemptive multitasking (managed by OS). & Yes (standard CPython blocks
CPU scaling). & I/O-bound tasks (waiting for sockets, files). & Medium
(thread stacks, context switches). & Shared memory (requires locks,
mutexes). \\
\textbf{Multiprocessing} (\passthrough{\lstinline!multiprocessing!}) &
Preemptive parallel execution (separate OS processes). & No (each
process has its own GIL/heap). & CPU-bound computations (math, data
parsing). & High (process fork/spawn cost, private heaps). & IPC (Pipes,
queues, shared memory, pickles). \\
\textbf{Asynchronous} (\passthrough{\lstinline!asyncio!} / coroutines) &
Cooperative multitasking (single thread, explicit yields). & Yes
(single-threaded). & High-concurrency network services (web servers). &
Low (coroutine objects are lightweight heap structures). & Direct
variable access (no locks required). \\
\bottomrule
\end{longtable}

\begin{center}\rule{0.5\linewidth}{0.5pt}\end{center}

\hypertarget{asyncio-event-loops-and-uvloop-internals}{%
\subsection{\texorpdfstring{27.2 Asyncio Event Loops and \texttt{uvloop}
Internals}{27.2 Asyncio Event Loops and uvloop Internals}}\label{asyncio-event-loops-and-uvloop-internals}}

Python's standard \passthrough{\lstinline!asyncio!} event loop uses
select-based system calls (\passthrough{\lstinline!selectors!} module,
e.g.~\passthrough{\lstinline!epoll!} on Linux,
\passthrough{\lstinline!kqueue!} on macOS) to track file descriptors.
While functional, the default implementation is written in pure Python
and introduces overhead when wrapping callback schedules.

\hypertarget{uvloop-optimization}{%
\subsubsection{1. uvloop Optimization}\label{uvloop-optimization}}

\passthrough{\lstinline!uvloop!} is a drop-in replacement for the
default asyncio event loop. Written in Cython and built on top of
\textbf{libuv} (the high-performance asynchronous I/O engine powering
Node.js), it replaces CPython's loop implementation with C-level
structures.

\begin{lstlisting}
CPython Default asyncio Loop:
[Asyncio Code] ---> [Pure Python Event Loop] ---> [selectors (epoll/kqueue)]

uvloop Event Loop:
[Asyncio Code] ---> [Cython Wrapper] ---> [libuv (C-native epoll/kqueue)]
\end{lstlisting}

Libuv optimizes execution paths by: * \textbf{System Call Reduction}:
Batching read/write operations to minimize transitions between user
space and kernel space. * \textbf{Direct Memory Buffers}: Allocating
internal buffers directly in C, avoiding intermediate Python byte
allocations. * \textbf{Zero-Overhead Timers}: Implementing timer queues
using binary heaps at the C level.

\hypertarget{configuring-uvloop-in-python}{%
\subsubsection{2. Configuring uvloop in
Python}\label{configuring-uvloop-in-python}}

To use \passthrough{\lstinline!uvloop!}, register it as the default
event loop policy at the entry point of your application:

\begin{lstlisting}[language=Python]
import asyncio
import sys

# Register uvloop policy on supported platforms
if sys.platform != 'win32':
    import uvloop
    asyncio.set_event_loop_policy(uvloop.EventLoopPolicy())

async def main():
    print("Running on optimized uvloop engine.")

if __name__ == '__main__':
    asyncio.run(main())
\end{lstlisting}

With \passthrough{\lstinline!uvloop!}, network throughput can increase
by 2x to 4x, reaching speeds comparable to implementations in Go or
Node.js.

\begin{center}\rule{0.5\linewidth}{0.5pt}\end{center}

\hypertarget{multiprocessing-shared-memory-ipc}{%
\subsection{27.3 Multiprocessing Shared Memory
IPC}\label{multiprocessing-shared-memory-ipc}}

Standard multiprocessing in Python relies on \textbf{Pipes} and
\textbf{Queues} for Inter-Process Communication (IPC). When a process
sends an object to another process: 1. The sender \textbf{pickles}
(serializes) the object into bytes. 2. The bytes are written to an OS
socket or pipe. 3. The receiver reads the bytes and \textbf{unpickles}
(deserializes) them to allocate new objects on its private heap.

For large data structures (like lists of floats or images), this
serialization pipeline degrades performance.

\hypertarget{shared-memory-architecture}{%
\subsubsection{1. Shared Memory
Architecture}\label{shared-memory-architecture}}

Python 3.8 introduced the
\passthrough{\lstinline!multiprocessing.shared\_memory!} module, which
maps a block of virtual memory across the address spaces of multiple OS
processes. Both processes can read and write to the same physical RAM
block directly, avoiding serialization overhead.

\begin{lstlisting}
Process 1 (Address Space)               Process 2 (Address Space)
  +--------------------+                  +--------------------+
  | Virtual Memory     |                  | Virtual Memory     |
  |  [Mapped Segment] -+----+        +----+-- [Mapped Segment] |
  +--------------------+    |        |    +--------------------+
                            v        v
                      +--------------------+
                      |    Physical RAM    |
                      |   [Shared Block]   |
                      +--------------------+
\end{lstlisting}

\hypertarget{robust-shared-memory-code-example}{%
\subsubsection{2. Robust Shared Memory Code
Example}\label{robust-shared-memory-code-example}}

Below is a complete script demonstrating parent and child processes
communicating using shared memory:

\begin{lstlisting}[language=Python]
import time
from multiprocessing import Process
from multiprocessing.shared_memory import SharedMemory

def child_process_worker(shm_name):
    # 1. Attach to existing shared memory block using unique name
    existing_shm = SharedMemory(name=shm_name)
    
    # 2. Access buffer as a memoryview array slice
    buffer = existing_shm.buf
    
    print(f"[Child] Connected to shared memory. Current contents: {bytes(buffer[:10])}")
    
    # Modify data in place (zero-copy modification)
    for i in range(10):
        buffer[i] = ord('A') + i
        
    print(f"[Child] Modified buffer contents in place.")
    
    # 3. Clean up the shared memory handle
    existing_shm.close()

def main():
    # 1. Allocate a shared memory block of 1024 bytes
    shm = SharedMemory(create=True, size=1024)
    shm_name = shm.name
    print(f"[Parent] Allocated Shared Memory block named: {shm_name}")
    
    # Initialize buffer contents
    shm.buf[:10] = b"0123456789"
    print(f"[Parent] Initial buffer contents: {bytes(shm.buf[:10])}")
    
    # 2. Spawn a child process, passing the shared memory block name
    p = Process(target=child_process_worker, args=(shm_name,))
    p.start()
    p.join()  # Wait for child process to finish writing
    
    # 3. Parent reads the updated data immediately without IPC serialization
    print(f"[Parent] Updated buffer contents read by Parent: {bytes(shm.buf[:10])}")
    
    # 4. Release and destroy shared memory block
    shm.close()
    shm.unlink()  # Instructs OS to free the shared memory resource

if __name__ == '__main__':
    main()
\end{lstlisting}

\begin{center}\rule{0.5\linewidth}{0.5pt}\end{center}
